\documentclass[10pt,a4paper]{amsart}
\usepackage[margin=19mm]{geometry}
\usepackage{amsmath,amssymb,mathtools}
\usepackage[scr=rsfs,cal=euler]{mathalfa}
\usepackage{xfrac}
\usepackage[unicode,hidelinks,bookmarks=false]{hyperref}
\newcommand{\ie}{i.e.\ }
\newcommand{\cf}{cf.\ }
\newcommand{\resp}{respectively\ }
\newcommand{\Subsection}{\subsection}
\providecommand{\nobreakdash}{\nobreak\hbox{-}}
\setlength{\emergencystretch}{2em}
\multlinegap0pt
\usepackage{amssymb}
\usepackage{algorithm}\usepackage{algpseudocode}
\usepackage{tikz}
\usepackage{mathtools}
\usepackage{xcolor}
\newtheorem{lemma}{Lemma}
\newcommand{\F}{\mathbb{F}}
\newcommand{\Fcal}{\mathcal{F}}
\newcommand{\Fix}{\mathrm{Fix}}
\newcommand{\GL}{\mathrm{GL}}
\newcommand{\Gam}{\Gamma}
\newcommand{\Hom}{\mathrm{Hom}}
\newcommand{\id}{\mathrm{id}}
\newcommand{\rank}{\operatorname{rank}}
\title{Almost All Vectorial Functions Have Trivial Extended-Affine Stabilizers}
\author{Keita Ishizuka}
\date{Source: arXiv:2602.06668v3 (CC BY 4.0)}
\begin{document}
\maketitle

\noindent\emph{Source and licence.} Almost All Vectorial Functions Have Trivial Extended-Affine Stabilizers. Version arXiv:2602.06668v3 (CC BY 4.0), distributed by arXiv under the Creative Commons Attribution 4.0 International licence, http://creativecommons.org/licenses/by/4.0/, as recorded in the arXiv licence metadata for this exact version and shown on the abstract page https://arxiv.org/abs/2602.06668v3. Full licence notice: \texttt{licenses/arxiv-2602.06668-CC-BY-4.0.txt}.\par\smallskip

\noindent\textit{Editorial scope.} Counted unit: the article's lemma lem:fix-g (source0.tex lines 369--377). The article's complete original proof of it is delivered with it (source0.tex lines 379--450, one \texttt{proof} environment of 72 lines). No mathematics is written, repaired or re-derived: the statement and the proof are the article's own lines, in the article's own notation, and no retained line is substituted or re-flowed.  Retained uncounted as the source-local context the proof needs: lines 332--346.  Nothing was omitted from any retained span, so the package carries no omission record.  No retained line contains a citation, so the wrapper carries no bibliography and no bibliography entry.  No retained line contains a figure environment, an image-inclusion command or drawing code, so no image is delivered and the package declares no asset dependency.  Editorial formatting only, in addition: an AMS \texttt{amsart} preamble that restates the article's own declarations and macros used by the retained lines, namely the environments \texttt{lemma} on separate counters, together with the standard packages those lines need.  The retained environments print in this document as Lemma 1 in order of appearance, whereas the article prints the same items with the numbering of its own full document.  The complete native source of the article as distributed in the arXiv e-print for this exact version is bundled unchanged as \texttt{sources/194114/source0.tex}.
\begin{lemma}[Orbit count]\label{lem:orbit-count}
Let $\sigma$ be a permutation of a finite set $X$ of size $N$.
A point $x\in X$ is a \emph{fixed point of $\sigma$} if $\sigma(x)=x$, and an
\emph{orbit of $\sigma$} is an orbit of the cyclic group $\langle\sigma\rangle$ acting
on $X$. Let $f$ denote the number of fixed points of $\sigma$ and $s$ the number of its
orbits. Then
\[
   s\le \frac{N+f}{2}.
\]
In particular, if $\sigma(x)=Px+a$ is a nontrivial affine permutation of $U=\F_q^n$,
then the number of orbits satisfies
\[
   s\le \frac{q^n+q^{n-1}}{2} = \frac{q+1}{2q}\cdot q^n.
\]
\end{lemma}

\begin{lemma}[Fixed-point bound for EA-group elements]\label{lem:fix-g}
Let $g=(P,a,Q,b,L)\in \Gam$ be nontrivial, where $P\in\GL(n,q)$, $a\in U$,
$Q\in\GL(m,q)$, $b\in W$, and $L\in\Hom(U,W)$.
Then the number of $F\in\Fcal$ satisfying $g\cdot F = F$ is bounded by
\[
   |\Fix(g)| \;\le\; q^{cmq^n}
\]
for some constant $c<1$ independent of $n,m$.
\end{lemma}

\begin{proof}
The fixed-point condition $g\cdot F=F$ reads
\begin{equation}\label{eq:fixed-eq}
   Q\,F(Px+a) + Lx + b = F(x)
   \quad(\forall x\in U).
\end{equation}
We distinguish two cases.

\medskip
\noindent\textbf{Case~1: $(P,a)\neq(I,0)$.}
Let $\sigma(x):=Px+a$; by assumption $\sigma\neq\id$.
Lemma~\ref{lem:orbit-count} shows that $\sigma$ has at most
\[
   s\le\frac{q+1}{2q}\cdot q^n
\]
orbits.
Decompose $U$ into $\sigma$-orbits:
\[
   U = \bigsqcup_{i=1}^s O_i,
   \qquad O_i = \{x_i,\sigma(x_i),\sigma^2(x_i),\dots,\sigma^{\ell_i-1}(x_i)\},
\]
where $\ell_i\ge 1$ is the length of orbit $O_i$.

Fix an orbit $O_i$ of length $\ell_i$.
Iterating~\eqref{eq:fixed-eq} starting from $x_i$ gives
\begin{align*}
   F(x_i)
      &= Q\,F(\sigma(x_i)) + Lx_i + b \\
      &= Q^2 F(\sigma^2(x_i)) + Q\bigl(L\sigma(x_i)+b\bigr) + Lx_i + b \\
      &\;\;\vdots \\
      &= Q^{\ell_i}F(x_i) + \sum_{j=0}^{\ell_i-1}Q^j\bigl(L\sigma^j(x_i)+b\bigr).
\end{align*}
Setting $r_i:=\sum_{j=0}^{\ell_i-1}Q^j\bigl(L\sigma^j(x_i)+b\bigr)\in W$ (which is
determined once $g$ and the orbit representative $x_i$ are fixed), this implies
\begin{equation}\label{eq:orbit-constraint}
   \bigl(I - Q^{\ell_i}\bigr) F(x_i) = r_i.
\end{equation}
The number of solutions $F(x_i)\in W$ to~\eqref{eq:orbit-constraint} is at most
$q^{\dim\ker(I-Q^{\ell_i})}\le q^m$.
Once $F(x_i)$ is chosen, the values $F(\sigma^k(x_i))$ for $1\le k<\ell_i$ are
uniquely determined by~\eqref{eq:fixed-eq}.

Repeating this argument for each orbit, we see that $F$ is determined by
choosing at most one free parameter (of size $\le q^m$) per orbit.
Since $s\le (q+1)/(2q)\cdot q^n$, the total number of admissible $F$ is at most
\[
   |\Fix(g)| \le (q^m)^s \;\le\; q^{m\cdot\frac{q+1}{2q}\cdot q^n}
   = q^{\frac{q+1}{2q}mq^n}.
\]
Hence $c=(q+1)/(2q)$ suffices in this case.

\medskip
\noindent\textbf{Case~2: $(P,a)=(I,0)$ but $(Q,b,L)\neq(I,0,0)$.}
Now~\eqref{eq:fixed-eq} becomes
\[
   Q\,F(x) + Lx + b = F(x)
   \quad\Longleftrightarrow\quad
   (I-Q)F(x) = Lx + b
   \qquad(\forall x\in U).
\]
If $Q\neq I$, then $\rank(I-Q)\ge 1$, so $\dim\ker(I-Q)\le m-1$, and for each
$x$ the equation $(I-Q)F(x)=Lx+b$ has at most $q^{m-1}$ solutions. Hence
\[
   |\Fix(g)| \;\le\; (q^{m-1})^{q^n} = q^{(m-1)q^n}.
\]
If $Q=I$, then the equation reduces to $Lx+b=0$ for all $x\in U$, which forces
$L=0$ and $b=0$, contradicting nontriviality. So no $F$ satisfies the equation,
and $|\Fix(g)|=0$.
In either subcase, $c = 1 - 1/m$ works when $m\ge 1$.

Combining both cases, we may take $c=\max\bigl\{(q+1)/(2q),1-1/m\bigr\}<1$.
\end{proof}

\end{document}
